\newcommand{\rev}{02 - 26-Mar-2021}
% REV02 Fri 26 Mar 2021 14:02:28 WIB
% REV01 Fri 26 Mar 2021 13:49:14 WIB
% START Fri 26 Mar 2021 12:10:02 WIB

\documentclass[12pt]{book}
\usepackage[a4paper, margin=50pt]{geometry}
\usepackage[pdftex]{graphicx}
\usepackage[hidelinks]{hyperref}
\usepackage{cellspace}
\usepackage{colortbl}
\usepackage{footnote}
\usepackage{lipsum}
\usepackage{multirow}
\usepackage{xcolor}
\setlength\cellspacetoplimit{4pt}
\setlength\cellspacebottomlimit{4pt}

% %%%%%

\newcommand{\pengarangs}{%
    Dickens McCartney\\
}
\newcommand{\judul}{%
Jenny Wren\\[13pt]
Our Mutual Friend
}

\begin{document}

\begin{titlepage}
    \begin{center}    

    \vspace*{15mm}
    \textbf{\Large \judul}

    \vspace*{30mm}       
    \textbf{by}

    \vspace*{15mm}    
    \textbf{\Large \pengarangs}

    \vspace*{4.0cm}

    \begin{center}
        \includegraphics[width=40mm]{ucls-coat}
    \end{center}

    \textbf{
       Universe Centra Le Sahara (UCLS)\\[11pt]
       Omar Bakry School of Management\\[11pt]
       Jabal Acacus Campus, Ghat. \\[11pt]
       Rev. \rev%
    }

    \vspace*{5mm}    
    \textbf{\LARGE \textcolor{red}{***** DRAFT *****}}

    \end{center}

\end{titlepage}

\pagenumbering{roman}

\tableofcontents

\newpage

\chapter*{Jenny Wren}
\addcontentsline{toc}{chapter}{Jenny Wren}

\begin{verbatim}
Like so many girls
Jenny Wren could sing
But a broken heart
Took her soul away

Like the other girls
Jenny Wren took wing
She could see the world
And it's foolish ways

How, we, spend our days
Casting, love aside
Loosing, site of life
Day, by, day

She saw poverty
Breaking all the home
Wounded warriors
Took her song away

But the day will come
Jenny Wren will sing
When this broken world
Mends its foolish ways

Now we, spend our days
Catching, up on life
All because of you
Jenny Wren

--- Dickens McCartney
\end{verbatim}

(Ghat, rev. \rev)

\newpage

\pagenumbering{arabic}

\chapter{ON THE LOOK OUT}

In these times of ours, though concerning the exact year there is no
need to be precise, a boat of dirty and disreputable appearance,
with two figures in it, floated on the Thames, between Southwark
bridge which is of iron, and London Bridge which is of stone, as an
autumn evening was closing in.

The figures in this boat were those of a strong man with ragged
grizzled hair and a sun-browned face, and a dark girl of nineteen or
twenty, sufficiently like him to be recognizable as his daughter.
The girl rowed, pulling a pair of sculls very easily; the man, with
the rudder-lines slack in his hands, and his hands loose in his
waistband, kept an eager look out. He had no net, hook, or line,
and he could not be a fisherman; his boat had no cushion for a
sitter, no paint, no inscription, no appliance beyond a rusty
boathook and a coil of rope, and he could not be a waterman; his
boat was too crazy and too small to take in cargo for delivery, and
he could not be a lighterman or river-carrier; there was no clue to
what he looked for, but he looked for something, with a most intent
and searching gaze. The tide, which had turned an hour before,
was running down, and his eyes watched every little race and eddy
in its broad sweep, as the boat made slight head-way against it, or
drove stern foremost before it, according as he directed his
daughter by a movement of his head. She watched his face as
earnestly as he watched the river. But, in the intensity of her look
there was a touch of dread or horror.

Allied to the bottom of the river rather than the surface, by reason
of the slime and ooze with which it was covered, and its sodden
state, this boat and the two figures in it obviously were doing
something that they often did, and were seeking what they often
sought. Half savage as the man showed, with no covering on his
matted head, with his brown arms bare to between the elbow and
the shoulder, with the loose knot of a looser kerchief lying low on
his bare breast in a wilderness of beard and whisker, with such
dress as he wore seeming to be made out of the mud that begrimed
his boat, still there was a business-like usage in his steady gaze.
So with every lithe action of the girl, with every turn of her wrist,
perhaps most of all with her look of dread or horror; they were
things of usage.

'Keep her out, Lizzie. Tide runs strong here. Keep her well afore
the sweep of it.'

Trusting to the girl's skill and making no use of the rudder, he eyed
the coming tide with an absorbed attention. So the girl eyed him.
But, it happened now, that a slant of light from the setting sun
glanced into the bottom of the boat, and, touching a rotten stain
there which bore some resemblance to the outline of a muffled
human form, coloured it as though with diluted blood. This caught
the girl's eye, and she shivered.

'What ails you?' said the man, immediately aware of it, though so
intent on the advancing waters; 'I see nothing afloat.'

The red light was gone, the shudder was gone, and his gaze, which
had come back to the boat for a moment, travelled away again.
Wheresoever the strong tide met with an impediment, his gaze
paused for an instant. At every mooring-chain and rope, at every
stationery boat or barge that split the current into a broad-
arrowhead, at the offsets from the piers of Southwark Bridge, at the
paddles of the river steamboats as they beat the filthy water, at the
floating logs of timber lashed together lying off certain wharves,
his shining eyes darted a hungry look. After a darkening hour or
so, suddenly the rudder-lines tightened in his hold, and he steered
hard towards the Surrey shore.





(TO BE CONTINUED)

% %%%%%%%%%%%%
\end{document}%
% End of document.
% %%%%%%%%%%%%%%%%


